\section{Complete tables of the calculations}
\label{app:tables}

Several sections of this paper quote selected numbers from calculations whose complete results
are tables. This appendix prints those tables, in the order of the sections to which they belong.
Each caption defines every column; the numbers are written by a script from the data files of the
program (Appendix~\ref{app:checks}). M and P stand throughout for the closures ``Momentum'' and
``Pressure'' of the perturbation equations (Sec.~\ref{sec:pert}), and $H_0$ is given in $\kmsMpc$.

\subsection{The paired state, the bound pair, temperature and vacuum (Secs.~\ref{sec:pair} and \ref{sec:cjt})}

Tables~\ref{tab:xthermo} and \ref{tab:xgl} follow the paired state from zero temperature to its
critical temperature at three moments of the deflation; they contain the universal ratios and the
coefficients of the Ginzburg-Landau energy that Sec.~\ref{sec:cjt} quotes.
Table~\ref{tab:xthreshold} compares two calculations of the critical value $D_c$ at which two
fermions bind. Table~\ref{tab:xthermal} lists the melting temperatures along the clock and
Table~\ref{tab:xvacuum} the fields at which the one-loop potential of the vacuum turns over.

\begin{table*}
\caption{\label{tab:xthermo}The paired state between zero temperature and the critical temperature
$T_c$, at three moments of the deflation, from the CJT potential, for $m_s=m$ and $\kF=0.5\,m$
(vector channel). $\lambda$: pairing strength. $\Delta_F(0)/T_c$: the gap at the Fermi momentum at zero
temperature divided by $T_c$; its weak-coupling value is 1.7639. Jump: the heat capacity of the paired
state minus that of the unpaired state, divided by the latter, at $T_c$; its weak-coupling value is
1.4261.}
\begin{ruledtabular}
\begin{tabular}{ccccccc}
$\td$ & $\gVS$ & $\lambda$ & $T_c/m$ & $T_c/\EF$ & $\Delta_F(0)/T_c$ & jump of the heat capacity \\
\hline
{{ROWS:tbl_cjt_thermo}}
\end{tabular}
\end{ruledtabular}
\end{table*}

\begin{table*}
\caption{\label{tab:xgl}The Ginzburg-Landau energy of the paired state at the three moments of
Table~\ref{tab:xthermo}. Near $T_c$ the CJT potential reduces to a quadratic term in the gap with the
coefficient $a(T)$, a quartic term with the coefficient $b$ and a gradient term with the coefficient
$K_{\mathrm{GL}}$. Columns 4 to 6: the slope of $a$ with respect to $\ln T$ at $T_c$, the coefficient
$b$ and the coefficient $K_{\mathrm{GL}}$, each divided by its value in the weak-coupling theory.
Last two columns: the coherence length of the Ginzburg-Landau energy, $(K_{\mathrm{GL}}/|a|)^{1/2}$ at
$|\ln(T/T_c)|=1$, in units of $1/m$ and in units of $v_F/T_c$ with the Fermi velocity $v_F$; in weak
coupling the latter is 0.1333.}
\begin{ruledtabular}
\begin{tabular}{cccccccc}
 & & & \multicolumn{3}{c}{ratio to weak coupling} & \multicolumn{2}{c}{coherence length} \\
$\td$ & $\gVS$ & $T_c/m$ & slope of $a(T)$ & $b$ & $K_{\mathrm{GL}}$ & $\times\,m$ & $\times\,T_c/v_F$ \\
\hline
{{ROWS:tbl_cjt_gl}}
\end{tabular}
\end{ruledtabular}
\end{table*}

\begin{table}
\caption{\label{tab:xthreshold}The critical value $D_c$ of $\DVS=\aVS\,m/m_s$ at which two fermions
bind (Sec.~\ref{sec:pair}), from the linearized gap equation in vacuum, whose kernel is relativistic,
and from the Schr\"odinger equation with the Yukawa potential. The two approach each other as the
radial mode becomes light and the binding non-relativistic.}
\begin{ruledtabular}
\begin{tabular}{cccc}
$m_s/m$ & gap equation & Schr\"odinger equation & difference \\
\hline
{{ROWS:tbl_threshold}}
\end{tabular}
\end{ruledtabular}
\end{table}

\begin{table*}
\caption{\label{tab:xthermal}Temperature along the clock: the melting of the expectation value of the
Vorton-Scalar, for three ratios $m_s/m$ and ten couplings each. $T_{\mathrm{high}}$: a closed estimate
of the temperature above which the expectation value vanishes. $T_{1/2}$ and $T_0$: the temperatures
at which the expectation value has fallen to one half and to 0.02 of its vacuum value, in the Hartree
approximation. Last column: ``no'' marks the moments at which no melting temperature can be given,
because the Hartree potential ceases to exist before the expectation value has fallen
(Sec.~\ref{sec:cjt}).}
\begin{ruledtabular}
\begin{tabular}{ccccccc}
$m_s/m$ & $\gVS$ & $C$ & $T_{\mathrm{high}}/m$ & $T_{1/2}/m$ & $T_0/m$ & melting defined \\
\hline
{{ROWS:tbl_thermal}}
\end{tabular}
\end{ruledtabular}
\end{table*}

\begin{table}
\caption{\label{tab:xvacuum}The one-loop potential of the Vorton-Scalar in the vacuum along the clock
($C\propto e^{\td}$): the field at which the potential has its maximum beyond the vacuum, and the
field beyond which it lies below its value in the vacuum, both in units of the vacuum value $\vVS$.}
\begin{ruledtabular}
\begin{tabular}{ccc}
$C$ & maximum & below the vacuum value \\
\hline
{{ROWS:tbl_vacuum}}
\end{tabular}
\end{ruledtabular}
\end{table}

\subsection{The unperturbed universe (Sec.~\ref{sec:cosmo})}

Table~\ref{tab:xdrift} gives the drift of the quantities of the model, Table~\ref{tab:xmassloss} the
examples of bound pairs that lose rest mass, and Table~\ref{tab:xgrowth} the growth of structure from
the two growth equations.
%%<1671>
Table~\ref{tab:xtemperature} gives the temperature of the quanta that the deflation creates at the
largest rates that the bounds allow.
%%</1671>

\begin{table*}
\caption{\label{tab:xdrift}How fast the quantities of the model drift. $n$: the logarithmic derivative
of the quantity with respect to the clock $\td$. Last three columns: the relative change
$n\,\gsum\,t_0$ of the quantity over the age $t_0$ of the universe at the largest total rate that each
bound of Table~\ref{tab:bounds} allows: lunar laser ranging (assumption A2); nucleosynthesis through
Newton's constant (A2); nucleosynthesis through the Fermi constant (assumption HV); with
$H_0t_0=0.9516$ of the model without deflation.}
\begin{ruledtabular}
\begin{tabular}{p{0.33\textwidth}p{0.2\textwidth}cccc}
 & & & lunar & \multicolumn{2}{c}{nucleosynthesis through} \\
quantity & where evaluated & $n$ & ranging & $\GN$ & $G_{\mathrm{Fermi}}$ \\
\hline
{{ROWS:tbl_drift}}
\end{tabular}
\end{ruledtabular}
\end{table*}

%%<1671>
\begin{table*}
\caption{\label{tab:xtemperature}The quanta that the deflation creates (Sec.~\ref{sec:modes}), for
three equal rates at the largest total rates that the bounds allow, and at $\gsum=H_0$. $\hbar\,\rt$:
the rate of one coordinate as an energy. $\Tdefl$: the temperature of Eq.~(\ref{eq:Tdefl}).
Exponent: $2\pi m/\rt$ for a particle of the mass $m$; the number of particles created per mode is
$e^{-\mathrm{exponent}}$.}
\begin{ruledtabular}
\begin{tabular}{lccccc}
bound & $\gsum/H_0$ & $\hbar\,\rt$ (eV) & $\Tdefl$ (K) & exponent, $m=1$~keV & exponent, $m=1$~GeV \\
\hline
{{ROWS:tbl_temperature}}
\end{tabular}
\end{ruledtabular}
\end{table*}
%%</1671>

\begin{table*}
\caption{\label{tab:xmassloss}Dark matter made of bound pairs, which lose rest mass as the binding
energy grows (Sec.~\ref{sec:cosmo}): four examples of $\DVS$, $m_s/m$ and $\aVS$ today, each at six
rates. ``Gravity feels the deflation'': yes, assumptions A1 and A2; no, Newton's constant is held
fixed and the deflation acts on the dark sector alone. Mass ratio: the rest mass of a pair today
divided by its rest mass in the early universe. $H_0$: the Hubble constant inferred with the density
of dark matter at early times, the density of baryons and the angle $\theta_*$ held fixed; change: its
difference from the value without deflation.}
\begin{ruledtabular}
\begin{tabular}{cccccccc}
$\DVS$ today & $m_s/m$ & $\aVS$ today & $\gsum/H_0$ & gravity feels the deflation & mass ratio & $H_0$ & change \\
\hline
{{ROWS:tbl_massloss}}
\end{tabular}
\end{ruledtabular}
\end{table*}

\begin{table*}
\caption{\label{tab:xgrowth}Linear growth of pressureless matter from $a_{\mathrm{sp}}=0.01$ to
today, from the growth equations (\ref{eq:growthM}) of the closure M and (\ref{eq:growthP}) of the
closure P. Columns 3 to 6, at the Hubble constant of column 2, which keeps the angle of the sound
horizon: the amplitude today and the growth rate $f_g$ today, each divided by its value without
deflation. Columns 7 and 8: the amplitude ratio at the fixed Hubble constant $67.36~\kmsMpc$. Last
column: the same from the equation
$d^{2}\delta_m/dx4^{2}+2H\,d\delta_m/dx4=4\pi\GN(x4)\rho_m\delta_m$, which is the equation of neither
closure and is shown for comparison.}
\begin{ruledtabular}
\begin{tabular}{ccccccccc}
 & & \multicolumn{2}{c}{amplitude} & \multicolumn{2}{c}{$f_g$} & \multicolumn{3}{c}{amplitude at fixed $H_0$} \\
$\gsum/H_0$ & $H_0$ & M & P & M & P & M & P & four-dimensional \\
\hline
{{ROWS:tbl_growth}}
\end{tabular}
\end{ruledtabular}
\end{table*}

\subsection{The microwave background (Secs.~\ref{sec:pert} and \ref{sec:spectra})}

Table~\ref{tab:xbackground} gives the thermal history for three models, Table~\ref{tab:xpeaks} the
acoustic peaks, Table~\ref{tab:xclosure} the difference between the two closures,
Table~\ref{tab:xgaps} what the agent must supply in each closure, and Tables~\ref{tab:xebgrowth} and
\ref{tab:xebgrowth2} the growth of structure in the complete calculation.

\begin{table*}
\caption{\label{tab:xbackground}Thermal history on the deflating spacetime (assumptions A1, A2; three
equal rates) for three models: the standard model (``fiducial''); the same with a fluid of extra
radiation of the density $N=0.5$; and the same with the stepped fluid of the model, $N=0.5$ and
$z_t=2\times10^{4}$. Held: the row keeps either the Hubble constant $67.36~\kmsMpc$ or the angle
$\theta_*$ of the same model without deflation. $z_*$: redshift of last scattering. $r_s$: sound
horizon at last scattering. $\YHe$: mass fraction of helium. $\GN$ ratio: Newton's constant at last
scattering divided by its value today.}
\begin{ruledtabular}
\begin{tabular}{lcccccccc}
model & $\gsum/H_0$ & held & $H_0$ & $z_*$ & $r_s$ (Mpc) & $100\,\theta_*$ & $\YHe$ & $\GN$ ratio \\
\hline
{{ROWS:tbl_eb_background}}
\end{tabular}
\end{ruledtabular}
\end{table*}

\begin{table*}
\caption{\label{tab:xpeaks}The first five acoustic peaks of the temperature spectrum at a fixed angle
of the sound horizon (closure M; without gravitational lensing): the multipole $\ell$ and the height
$D_\ell^{TT}$ in $\mu\mathrm{K}^{2}$, without deflation and at the largest rates that lunar laser
ranging and nucleosynthesis allow.}
\begin{ruledtabular}
\begin{tabular}{ccccccc}
 & \multicolumn{2}{c}{no deflation} & \multicolumn{2}{c}{lunar ranging rate} & \multicolumn{2}{c}{nucleosynthesis rate} \\
peak & $\ell$ & $D_\ell^{TT}$ & $\ell$ & $D_\ell^{TT}$ & $\ell$ & $D_\ell^{TT}$ \\
\hline
{{ROWS:tbl_eb_peaks}}
\end{tabular}
\end{ruledtabular}
\end{table*}

\begin{table*}
\caption{\label{tab:xclosure}The two closures compared, at a fixed angle of the sound horizon.
Difference: the largest difference between the temperature spectra of the two closures in the range
of multipoles named, in units of the cosmic variance of a single multipole. Share of the effect: the
root of the sum of the squares of that difference over the range, divided by the same for the effect
of the deflation itself. $\sigma_8$ ratio: the amplitude of clustering divided by its value without
deflation. For $\ell\ge30$ the closure does not matter; for $\ell<30$ the closures differ by more
than the effect of the deflation in either.}
\begin{ruledtabular}
\begin{tabular}{ccccccc}
 & \multicolumn{2}{c}{$\ell\ge30$} & \multicolumn{2}{c}{$\ell<30$} & \multicolumn{2}{c}{$\sigma_8$ ratio} \\
$\gsum/H_0$ & difference & share of the effect & difference & share of the effect & M & P \\
\hline
{{ROWS:tbl_eb_closure}}
\end{tabular}
\end{ruledtabular}
\end{table*}

\begin{table*}
\caption{\label{tab:xgaps}What the agent must supply in each closure, at the largest rate that
nucleosynthesis allows, for the wavenumber $k=0.02~\mathrm{Mpc}^{-1}$ at five conformal times.
$\gsum/H$: the ratio $a_{\mathrm{sp}}\gsum/\Hcf$. $\mathcal{U}=\etasy+\alsy'$: the sum of the two
gravitational potentials of the conformal Newtonian gauge, the potential that light feels. The gaps
$\cYM$, $\cYT$, $\cYS$ of Eq.~(\ref{eq:gaps}) are divided by $\gsum/H$ times the size of the terms of
their equations, $k^{2}\Hcf|\mathcal{U}|$, $(k^{2}+\Hcf^{2})|\mathcal{U}|$ and $|\mathcal{U}|$, so
that a number of order one means ``of the first order in the rate''. Last two columns: the residuals
of the identities (\ref{eq:id2}) and (\ref{eq:id1}), relative to the size of their terms.}
\begin{ruledtabular}
\begin{tabular}{ccccccccc}
closure & $\tcf$ (Mpc) & $\gsum/H$ & $\mathcal{U}$ & $\cYM$ & $\cYT$ & $\cYS$ & residual of (\ref{eq:id2}) & residual of (\ref{eq:id1}) \\
\hline
{{ROWS:tbl_eb_gaps}}
\end{tabular}
\end{ruledtabular}
\end{table*}

\begin{table}
\caption{\label{tab:xebgrowth}Growth of the density contrast of cold dark matter and baryons at
$k=0.1~\mathrm{Mpc}^{-1}$ from $a_{\mathrm{sp}}=0.01$ to today, divided by its value without
deflation, at the fixed Hubble constant $67.36~\kmsMpc$: from the complete Einstein-Boltzmann
calculation (``program'') and from the growth equation of the same closure,
Eq.~(\ref{eq:growthM}) for M and Eq.~(\ref{eq:growthP}) for P.}
\begin{ruledtabular}
\begin{tabular}{ccccc}
 & \multicolumn{2}{c}{closure M} & \multicolumn{2}{c}{closure P} \\
$\gsum/H_0$ & program & equation & program & equation \\
\hline
{{ROWS:tbl_eb_growth}}
\end{tabular}
\end{ruledtabular}
\end{table}

\begin{table}
\caption{\label{tab:xebgrowth2}The same growth in the complete calculation at the Hubble
constant that keeps the angle of the sound horizon, and the amplitude $\sigma_8$ divided by its value
without deflation.}
\begin{ruledtabular}
\begin{tabular}{cccccc}
 & & \multicolumn{2}{c}{growth} & \multicolumn{2}{c}{$\sigma_8$ ratio} \\
$\gsum/H_0$ & $H_0$ & M & P & M & P \\
\hline
{{ROWS:tbl_eb_growth_angle}}
\end{tabular}
\end{ruledtabular}
\end{table}

\subsection{Fits to mock data (Sec.~\ref{sec:forecast})}

Table~\ref{tab:xfcprofile} lists the best fits of the six parameters of the standard model at fixed
rates of deflation. Tables~\ref{tab:xfamfisher} to \ref{tab:xfamjoint} give, for the seven families
of extra radiation, the Fisher forecasts, the best fits at fixed density $N$ of the extra radiation,
their summary, and the forecasts with $N$ and the rate free together. ``Model fluid'' is the stepped
fluid of the model with the transition redshift $z_t$; ``with deep survey'' means the Planck-like
mock data together with a deep survey of a quarter of the sky.

\begin{table*}
\caption{\label{tab:xfcprofile}Best fits of the six parameters of the standard model to mock data of
the microwave background at fixed total rates of deflation (three equal rates; assumptions A1, A2,
A3). $\Delta\chi^{2}$: relative to the best fit without deflation, which is zero by construction.
$H_0$, $\omega_c$, $n_s$: the best-fitting Hubble constant, density of cold dark matter and spectral
index. $\sigma_8$: the amplitude of clustering of the best fit. $\GN$ ratio: Newton's constant at last
scattering divided by its value today.}
\begin{ruledtabular}
\begin{tabular}{llccccccc}
mock data & closure & $\gsum/H_0$ & $\Delta\chi^{2}$ & $H_0$ & $\omega_c$ & $n_s$ & $\sigma_8$ & $\GN$ ratio \\
\hline
{{ROWS:tbl_fc_profile}}
\end{tabular}
\end{ruledtabular}
\end{table*}

\begin{table*}
\caption{\label{tab:xfamfisher}Fisher forecasts for the seven families of extra radiation, with static
extra coordinates (``static'') and on the deflating spacetime at the largest rate that lunar laser
ranging allows (``deflating''): the standard error of the density $N$ with the six parameters of the
standard model free, and the change of the best-fit Hubble constant per unit of $N$.}
\begin{ruledtabular}
\begin{tabular}{llcccc}
 & & \multicolumn{2}{c}{error of $N$} & \multicolumn{2}{c}{$dH_0/dN$} \\
mock data & family & static & deflating & static & deflating \\
\hline
{{ROWS:tbl_fam_fisher}}
\end{tabular}
\end{ruledtabular}
\end{table*}

\begin{table*}
\caption{\label{tab:xfamprofiles}Best fits of the six parameters of the standard model at fixed
density $N$ of extra radiation, for the seven families and both sets of mock data, with static extra
coordinates and on the deflating spacetime at the largest rate that lunar laser ranging allows. $H_0$:
the best-fitting Hubble constant; shift: the difference of the two. $\Delta\chi^{2}$: relative to the
best fit of the standard model with static extra coordinates.}
\begin{ruledtabular}
\begin{tabular}{llcccccc}
 & & & \multicolumn{3}{c}{$H_0$} & \multicolumn{2}{c}{$\Delta\chi^{2}$} \\
mock data & family & $N$ & static & deflating & shift & static & deflating \\
\hline
{{ROWS:tbl_fam_profiles}}
\end{tabular}
\end{ruledtabular}
\end{table*}

\begin{table*}
\caption{\label{tab:xfamsummary}Summary of the fits of Table~\ref{tab:xfamprofiles}, by interpolation
in $N$: the density $N$ at which the best fit reaches $\Delta\chi^{2}=4$; the Hubble constant there;
and the $\Delta\chi^{2}$ of the best fit whose Hubble constant equals the local measurement
$73.50~\kmsMpc$. An asterisk marks a value obtained by extrapolation beyond the largest $N$ computed.}
\begin{ruledtabular}
\begin{tabular}{lcccccc}
 & \multicolumn{2}{c}{$N$ at $\Delta\chi^{2}=4$} & \multicolumn{2}{c}{$H_0$ there} & \multicolumn{2}{c}{$\Delta\chi^{2}$ at $H_0=73.50$} \\
family & static & deflating & static & deflating & static & deflating \\
\hline
\multicolumn{7}{c}{Planck-like mock data} \\
{{ROWS:tbl_fam_summary_planck}}
\hline
\multicolumn{7}{c}{with deep survey} \\
{{ROWS:tbl_fam_summary_deep}}
\end{tabular}
\end{ruledtabular}
\end{table*}

\begin{table*}
\caption{\label{tab:xfamjoint}The density $N$ of extra radiation and the total rate of deflation as
free parameters together (Fisher matrices with eight parameters at the standard model). Columns 2 and
3: the standard error of $N$ with the rate held at zero and with the rate free. Column 4: the standard
error of $\gsum/H_0$ with $N$ free. Correlation: the correlation coefficient of $N$ and the rate. $N$
per unit rate: the change of the best-fitting $N$ per unit of $\gsum/H_0$. Net slope: the change of
the best-fit Hubble constant per unit of $\gsum/H_0$ when $N$ follows its best-fitting value. Largest
rise: the largest rise of the best-fit Hubble constant among the models with $N\ge0$, $\gsum\ge0$ and
$\Delta\chi^{2}\le3.84$, which is reached at zero rate.}
\begin{ruledtabular}
\begin{tabular}{lccccccc}
 & \multicolumn{2}{c}{error of $N$} & error of & & $N$ per & net & largest \\
family & rate zero & rate free & the rate & correlation & unit rate & slope & rise \\
\hline
\multicolumn{8}{c}{Planck-like mock data} \\
{{ROWS:tbl_fam_joint_planck}}
\hline
\multicolumn{8}{c}{with deep survey} \\
{{ROWS:tbl_fam_joint_deep}}
\end{tabular}
\end{ruledtabular}
\end{table*}

\subsection{Superfluid dark matter (Sec.~\ref{sec:dm})}

Table~\ref{tab:xsfdmhalf} lists the wavenumber of the half mode for the fifteen models of
Sec.~\ref{sec:dm} and Table~\ref{tab:xsfdmbounds} the limits that follow from the Lyman-$\alpha$
forest.

\begin{table*}
\caption{\label{tab:xsfdmhalf}Superfluid dark matter: the wavenumber of the half mode, at which the
power spectrum of matter first falls to one quarter of that of cold dark matter, in
$\mathrm{Mpc}^{-1}$, for a degenerate gas of the fermions of eight masses and for a condensate of
bound pairs with seven sound speeds $c_{s0}$ today: without deflation, at the largest rate that lunar
laser ranging allows, and at the largest rate that nucleosynthesis allows in the two closures. Last
column: the change at the rate of nucleosynthesis relative to no deflation, in per cent.}
\begin{ruledtabular}
\begin{tabular}{lccccc}
 & no & lunar ranging & \multicolumn{2}{c}{nucleosynthesis rate} & change \\
dark matter & deflation & rate & M & P & (per cent) \\
\hline
{{ROWS:tbl_sfdm_half}}
\end{tabular}
\end{ruledtabular}
\end{table*}

\begin{table*}
\caption{\label{tab:xsfdmbounds}Limits on superfluid dark matter from the requirement that its half
mode lies at a wavenumber not smaller than that of thermal warm dark matter at two mass limits of the
Lyman-$\alpha$ forest (``relic'', the mass of the thermal relic). Half mode: the wavenumber required.
Fermion mass: the smallest mass of the fermions of a degenerate gas that meets it. Largest $c_{s0}/c$:
the largest sound speed today of a condensate of bound pairs that meets it, also in m/s.
$R_{\mathrm{TF}}$: the radius of a core of the condensate supported by its own pressure, at that
sound speed.}
\begin{ruledtabular}
\begin{tabular}{lcccccc}
deflation & relic (keV) & half mode ($\mathrm{Mpc}^{-1}$) & fermion mass (keV) & largest $c_{s0}/c$ & in m/s & $R_{\mathrm{TF}}$ (pc) \\
\hline
{{ROWS:tbl_sfdm_bounds}}
\end{tabular}
\end{ruledtabular}
\end{table*}

\subsection{Symmetries (Secs.~\ref{sec:geometry} and \ref{sec:hidden})}

Table~\ref{tab:xisometries} counts the rotations and boosts that remain symmetries of the deflating
spacetime, and Table~\ref{tab:xsymmetry} lists, item by item, the internal symmetry of the fields that
depend on $x1,x2,x3,x4$ only, as the program computed it for the deflating spacetime.

\begin{table*}
\caption{\label{tab:xisometries}Of the 28 rotations and boosts of flat spacetime, the number that
remain symmetries of the deflating spacetime for three sets of rates
$(\gff,\gss,\gvv)$; and of the six rotations of the extra axes $x0,x5,x6,x7$, whose generators act as
internal symmetries on fields that do not depend on these coordinates, the number that remain
symmetries of the spacetime.}
\begin{ruledtabular}
\begin{tabular}{lcccccccc}
 & rotations & \multicolumn{3}{c}{symmetries for the rates} & internal & \multicolumn{3}{c}{of these, symmetries for the rates} \\
spacetime & and boosts & $(1,2,3)$ & $(2,2,5)$ & $(1,1,1)$ & generators & $(1,2,3)$ & $(2,2,5)$ & $(1,1,1)$ \\
\hline
{{ROWS:tbl_isometries}}
\end{tabular}
\end{ruledtabular}
\end{table*}

\begin{table*}
\caption{\label{tab:xsymmetry}The internal symmetry of the fields that depend on $x1,x2,x3,x4$ only
(Sec.~\ref{sec:hidden}), item by item, computed with the matrices of the \SIG{$4+4$}{$7+1$} spacetime;
a number of the order of $10^{-16}$ is zero to the accuracy of the arithmetic. $P$: the traceless
generator of $U(1)_E$. $T_i$: the generators of one of the two ideals $su(2)_A$, $su(2)_B$. The pair
operators are the four vector pair operators of Eq.~(\ref{eq:pairs}).}
\begin{ruledtabular}
\begin{tabular}{p{0.46\textwidth}p{0.16\textwidth}p{0.3\textwidth}}
item & value & expected, and what it means \\
\hline
{{ROWS:tbl_symmetry_ms\SIG{|cols=0,2,3}{|cols=0,1,3}}}
\end{tabular}
\end{ruledtabular}
\end{table*}

% all tables of this appendix are printed before the next appendix begins
\clearpage
